\begin{abstract}
When a Linux kernel API changes, hundreds of drivers are converted by hand over months, and the earlier
conversions are reviewed examples for the later ones. We ask whether small open language models running on a
laptop can turn such examples into Coccinelle semantic patches, the rules kernel maintainers already use for
tree-wide changes. We adapt guided code generation to this task: the training examples are decomposed by script
into a tree of change variants, the model writes one small rule per variant under feedback from applying it to the
examples, deterministic components repair SmPL syntax and generalise driver-specific names, rules are composed
greedily, and a compile critic repairs or rejects a rule whose output does not build. We evaluate on a
pre-registered test set of \testMigrations{} migrations whose new APIs appeared after the training cutoffs of all
three models (Qwen2.5-Coder 7B, Llama~3.1 8B, Gemma~3 4B), with \testBuildable{} held-out driver conversions,
counting a conversion as solved only if the rule's output matches the human fix and compiles. With the framework,
Qwen2.5-Coder solves \testQwenCoderFull{} conversions on average over three seeds (range
\testQwenCoderFullMin--\testQwenCoderFullMax), against \testQwenCoderOneShot{} when prompted directly; Spinfer
solves \testSpinfer{} and a frontier LLM agent \testFrontier. The local model solves fewer conversions than
Spinfer but its rules are the most precise on the test set (\testSafeQwenCoderPrecision{} of the files they change
are solved, against \testSafeSpinferPrecision{} for Spinfer and \testSafeFrontierPrecision{} for the frontier
agent), and the compile critic removes every output that matches the human fix but does not build. The two general-purpose
models stay weak even with guidance (at most \testLlamaFullMax{} solved). All code, prompts, model responses and
results are released.
\end{abstract}
